\section{Euclidean Reference for Self-Study}

The following material preserves the geometric background behind the Cartesian construction. It is included so that the chapter remains self-contained, but it is not intended to occupy the same lecture time as coordinates, graphs, and intersections.

\subsection{Three Different Levels Must Not Be Confused}

\begin{corebox}[The logical hierarchy]
\begin{enumerate}
    \item \emph{Euclid's postulates} describe the geometric world being assumed.
    \item \emph{Straightedge-and-compass rules} describe which elementary
          drawings and intersections may be used in a construction.
    \item \emph{Derived constructions} must be built from those elementary
          operations and justified by a proof.
\end{enumerate}
\end{corebox}

\begin{remarkbox}
A perpendicular bisector, a perpendicular through a point, an angle bisector,
or a parallel through a point is not treated as a primitive permission. Each
must be constructed from the elementary toolkit and then justified.
\end{remarkbox}

\subsection{Euclid's Five Postulates}

The following are the five postulates used in the first book of Euclid's
\emph{Elements}. They specify the elementary geometric operations and relations
from which the later constructions begin. The diagrams illustrate their meaning;
they are not proofs.

\begin{corebox}[Postulate I: joining two points]
A straight segment can be drawn from any point to any other point.
\end{corebox}

\begin{center}
\begin{tikzpicture}[scale=0.9]
    \coordinate (A) at (-2.4,-0.6);
    \coordinate (B) at (2.1,0.8);
    \draw[subset] (A) -- (B);
    \fill[geometricSubsetColor] (A) circle (2pt) node[below] {$A$};
    \fill[geometricSubsetColor] (B) circle (2pt) node[above] {$B$};
    \node at (0,-1.55) {a segment is drawn from $A$ to $B$};
\end{tikzpicture}
\end{center}

\begin{corebox}[Postulate II: extending a segment]
A finite straight segment can be extended continuously in the same straight
line.
\end{corebox}

\begin{center}
\begin{tikzpicture}[scale=0.9]
    \coordinate (A) at (-2.2,0);
    \coordinate (B) at (1.2,0);
    \coordinate (C) at (3.1,0);
    \draw[subset] (A) -- (B);
    \draw[helper,->] (B) -- (C);
    \fill[geometricSubsetColor] (A) circle (2pt) node[below] {$A$};
    \fill[geometricSubsetColor] (B) circle (2pt) node[below] {$B$};
    \node[below] at (C) {$C$};
    \node at (0,-1.15) {$AB$ is extended beyond $B$};
\end{tikzpicture}
\end{center}

\begin{corebox}[Postulate III: drawing a circle]
A circle can be drawn with any chosen centre and any chosen radius.
\end{corebox}

\begin{center}
\begin{tikzpicture}[scale=0.9]
    \coordinate (O) at (0,0);
    \coordinate (A) at (2.2,0);
    \draw[helper] (O) circle (2.2);
    \draw[subset] (O) -- (A) node[midway,below] {$r$};
    \fill[geometricSubsetColor] (O) circle (2pt) node[below left] {$O$};
    \fill[geometricSubsetColor] (A) circle (2pt) node[right] {$A$};
    \node at (0,-2.75) {centre $O$ and radius $r$ determine a circle};
\end{tikzpicture}
\end{center}

\begin{interpretationbox}[Why the circle postulate matters]
This postulate permits a length to be copied geometrically without first
assigning a number to it. The compass transfers a geometric magnitude directly.
\end{interpretationbox}

\begin{corebox}[Postulate IV: equality of right angles]
All right angles are congruent.
\end{corebox}

\begin{center}
\begin{tikzpicture}[scale=0.85]
    \draw[subset] (-3.2,0) -- (-0.2,0);
    \draw[subset] (-1.7,-1.3) -- (-1.7,1.6);
    \draw (-1.42,0) -- (-1.42,0.28) -- (-1.7,0.28);

    \draw[subset] (0.8,-0.8) -- (3.5,0.7);
    \draw[subset] (2.9,-1.25) -- (1.4,1.45);
    \draw (2.24,0.02) -- (2.10,0.27) -- (1.85,0.13);

    \node at (0,-2.0) {every right angle represents the same geometric magnitude};
\end{tikzpicture}
\end{center}

\begin{remarkbox}
This postulate does not say that a perpendicular may simply be drawn without
construction. It says that once right angles have been constructed, they are all
congruent.
\end{remarkbox}

\begin{corebox}[Postulate V: the parallel postulate]
If a transversal meets two straight lines and the interior angles on one side
sum to less than two right angles, then the two lines, when extended, meet on
that side.
\end{corebox}

\begin{center}
\begin{tikzpicture}[scale=0.85]
    \draw[subset] (-3.4,1.15) -- (3.5,0.15) node[right] {$\ell$};
    \draw[subset] (-3.4,-1.2) -- (3.5,-0.35) node[right] {$m$};
    \draw[helper] (-1.0,-2.0) -- (0.6,2.0) node[above] {$t$};
    \node at (0,-2.45) {the angle condition determines on which side $\ell$ and $m$ meet};
\end{tikzpicture}
\end{center}

\begin{interpretationbox}[Modern working form]
A familiar equivalent is Playfair's formulation: through a point outside a
given line there passes exactly one parallel to that line. We shall use that
form later, after the relation between the two formulations has been made clear.
\end{interpretationbox}

\begin{summarybox}[What the postulates provide]
The first three postulates provide elementary constructions, the fourth makes
right angles comparable, and the fifth controls parallelism. Together they
specify the geometric world available before coordinates and algebra are
introduced.
\end{summarybox}

\subsection{Elementary Rules for Straightedge-and-Compass Work}

For explicit constructions we use the following operational rules.

\begin{enumerate}
    \item Through two already constructed distinct points, draw the straight line.
    \item With an already constructed point as centre and an already constructed
          segment as radius, draw a circle.
    \item Retain intersection points of two constructed lines, a line and a
          circle, or two circles whenever such intersections exist.
\end{enumerate}

\begin{center}
\begin{tikzpicture}[scale=0.9]
    \coordinate (A) at (-2.3,-0.4);
    \coordinate (B) at (2.1,0.9);
    \draw[subset] (A) -- (B);
    \fill[geometricSubsetColor] (A) circle (2pt) node[below] {$A$};
    \fill[geometricSubsetColor] (B) circle (2pt) node[above] {$B$};
    \node at (0,-1.35) {draw the line through two constructed points};
\end{tikzpicture}
\end{center}

\begin{center}
\begin{tikzpicture}[scale=0.9]
    \coordinate (O) at (0,0);
    \coordinate (R) at (2.0,0);
    \draw[helper] (O) circle (2.0);
    \draw[subset] (O) -- (R) node[midway,below] {$r$};
    \fill[geometricSubsetColor] (O) circle (2pt) node[below left] {$O$};
    \fill[geometricSubsetColor] (R) circle (2pt);
    \node at (0,-2.55) {draw a circle from a constructed centre and radius};
\end{tikzpicture}
\end{center}

\begin{center}
\begin{tikzpicture}[scale=0.85]
    \coordinate (C) at (-1.5,0);
    \coordinate (D) at (1.5,0);
    \draw[helper] (C) circle (2.15);
    \draw[helper] (D) circle (2.15);
    \fill (C) circle (1.7pt) node[left] {$C$};
    \fill (D) circle (1.7pt) node[right] {$D$};
    \fill[geometricSubsetColor] (0,1.55) circle (2pt) node[above] {$P$};
    \fill[geometricSubsetColor] (0,-1.55) circle (2pt) node[below] {$Q$};
    \node at (0,-2.65) {retain the intersection points as new constructible points};
\end{tikzpicture}
\end{center}

These rules are the actual construction toolkit. The constructions below must
be reduced to repeated uses of them.

\subsection{Constructions Derived from the Toolkit}

The following objects are not added as new primitive permissions. Each is
constructed from lines, circles, and their intersections.

\subsubsection{The Perpendicular Bisector and the Midpoint}

\begin{derivationbox}[Construction]
Let $A$ and $B$ be distinct points. Draw two circles with the same radius,
one centred at $A$ and one centred at $B$, choosing the radius larger than half
of $|AB|$. Let their intersection points be $P$ and $Q$.

Because $P$ lies on both circles,
\[
|PA|=|PB|,
\]
and similarly,
\[
|QA|=|QB|.
\]
Thus both $P$ and $Q$ are equidistant from $A$ and $B$.
\end{derivationbox}

\begin{center}
\begin{tikzpicture}[scale=0.85]
    \coordinate (A) at (-2.2,0);
    \coordinate (B) at (2.2,0);
    \coordinate (P) at (0,2.4);
    \coordinate (Q) at (0,-2.4);
    \draw[helper] (A) circle (3.25);
    \draw[helper] (B) circle (3.25);
    \draw[subset] (A) -- (B);
    \draw[very thick] (0,-3.0) -- (0,3.0);
    \fill[geometricSubsetColor] (A) circle (2pt) node[below] {$A$};
    \fill[geometricSubsetColor] (B) circle (2pt) node[below] {$B$};
    \fill[geometricSubsetColor] (P) circle (2pt) node[above] {$P$};
    \fill[geometricSubsetColor] (Q) circle (2pt) node[below] {$Q$};
\end{tikzpicture}
\end{center}

\begin{theorembox}
The line $PQ$ is the perpendicular bisector of $AB$. Its intersection with
$AB$ is the midpoint of $AB$. Thus both midpoint and perpendicularity are
constructed from equal radii and intersections.
\end{theorembox}

\subsubsection{A Perpendicular Through a Point on a Line}

\begin{derivationbox}[Construction]
Let $P$ lie on a line $\ell$.

\textbf{Step 1.} Draw a circle centred at $P$. It meets $\ell$ at two points
$A$ and $B$. Since $|PA|=|PB|$, the point $P$ is the midpoint of $AB$.

\textbf{Step 2.} Apply the perpendicular-bisector construction to $AB$. The
resulting perpendicular bisector passes through its midpoint $P$.
\end{derivationbox}

\begin{center}
\begin{tikzpicture}[scale=0.9]
    \coordinate (A) at (-2.2,0);
    \coordinate (P) at (0,0);
    \coordinate (B) at (2.2,0);
    \coordinate (Q) at (0,2.4);
    \coordinate (R) at (0,-2.4);

    \draw[subset] (-3.5,0) -- (3.5,0) node[right] {$\ell$};
    \draw[dashed,line width=0.9pt,draw=blue!65] (P) circle (2.2);
    \draw[dashed,line width=0.9pt,draw=orange!85!black] (A) circle (3.25);
    \draw[dashed,line width=0.9pt,draw=orange!85!black] (B) circle (3.25);
    \draw[very thick] (Q) -- (R) node[below] {$m=QR$};
    \draw (0.28,0) -- (0.28,0.28) -- (0,0.28);

    \fill[geometricSubsetColor] (A) circle (2pt) node[below] {$A$};
    \fill[geometricSubsetColor] (P) circle (2pt) node[below right] {$P$};
    \fill[geometricSubsetColor] (B) circle (2pt) node[below] {$B$};
    \fill[geometricSubsetColor] (Q) circle (2pt) node[above] {$Q$};
    \fill[geometricSubsetColor] (R) circle (2pt) node[below] {$R$};

    \node[blue!65,align=left] at (-3.0,2.75)
        {\small Step 1: circle centred at $P$};
    \node[orange!85!black,align=left] at (3.0,-2.75)
        {\small Step 2: circles centred at $A$ and $B$};
\end{tikzpicture}
\end{center}

\begin{theorembox}
The constructed line $m$ passes through $P$ and is perpendicular to $\ell$.
\end{theorembox}

\subsubsection{A Perpendicular Through a Point Outside a Line}

\begin{derivationbox}[Construction]
Let $P$ lie outside a line $\ell$.

\textbf{Step 1.} Draw a circle centred at $P$ with radius large enough to meet
$\ell$ at two points $A$ and $B$. Then $|PA|=|PB|$, so $P$ lies on the
perpendicular bisector of $AB$.

\textbf{Step 2.} Draw equal-radius circles centred at $A$ and $B$, chosen so that
they meet again at a second point $Q$. The points $P$ and $Q$ are both
equidistant from $A$ and $B$.
\end{derivationbox}

\begin{center}
\begin{tikzpicture}[scale=0.9]
    \coordinate (P) at (0,2.0);
    \coordinate (Q) at (0,-2.0);
    \coordinate (A) at (-2.1,0);
    \coordinate (B) at (2.1,0);

    \draw[subset] (-3.5,0) -- (3.5,0) node[right] {$\ell$};
    \draw[dashed,line width=0.9pt,draw=blue!65] (P) circle (2.9);
    \draw[dashed,line width=0.9pt,draw=orange!85!black] (A) circle (2.9);
    \draw[dashed,line width=0.9pt,draw=orange!85!black] (B) circle (2.9);
    \draw[very thick] (0,-2.8) -- (0,2.8) node[above] {$m=PQ$};
    \draw (0.28,0) -- (0.28,0.28) -- (0,0.28);

    \fill[geometricSubsetColor] (P) circle (2pt) node[right] {$P$};
    \fill[geometricSubsetColor] (Q) circle (2pt) node[right] {$Q$};
    \fill[geometricSubsetColor] (A) circle (2pt) node[below] {$A$};
    \fill[geometricSubsetColor] (B) circle (2pt) node[below] {$B$};

    \node[blue!65,align=left] at (-3.0,2.65)
        {\small Step 1: circle centred at $P$};
    \node[orange!85!black,align=left] at (3.0,-2.65)
        {\small Step 2: circles centred at $A$ and $B$};
\end{tikzpicture}
\end{center}

\begin{theorembox}
The line $PQ$ is the perpendicular bisector of $AB$ and therefore is
perpendicular to $\ell$ through the prescribed external point $P$.
\end{theorembox}

\subsubsection{Bisecting an Angle}

\begin{derivationbox}[Construction and justification]
Let two rays with common endpoint $O$ form an angle. Draw a circle centred at
$O$ meeting the rays at $A$ and $B$. Then draw equal-radius circles centred at
$A$ and $B$, and let one of their intersections inside the angle be $P$. Since
\[
|OA|=|OB|,
\qquad
|AP|=|BP|,
\qquad
|OP|=|OP|,
\]
the triangles $OAP$ and $OBP$ are congruent.
\end{derivationbox}

\begin{theorembox}
The congruence of $OAP$ and $OBP$ gives
\[
\angle AOP=\angle POB.
\]
Therefore the ray $OP$ bisects the angle.
\end{theorembox}

\subsubsection{Constructing a Parallel Through a Point}

\begin{derivationbox}[Construction]
Let $P$ lie outside a line $\ell$. First construct a line $m$ through $P$
perpendicular to $\ell$. Then construct a line $n$ through $P$ perpendicular
to $m$.

The transversal $m$ forms equal corresponding right angles with $n$ and
$\ell$. By the converse of the corresponding-angles theorem, the two lines are
parallel.
\end{derivationbox}

\begin{theorembox}
The constructed line satisfies
\[
\boxed{n\parallel\ell}.
\]
The fifth postulate, equivalently Playfair's postulate, guarantees that this is
the unique parallel to $\ell$ through $P$.
\end{theorembox}

\subsection{Parallel Lines and Proportional Segments}

A central fact of Euclidean geometry is that parallel lines preserve ratios.
This is the content of Thales' theorem in the form needed here.

\begin{theorembox}
Let $D$ lie on $AB$ and $E$ lie on $AC$ in a triangle $ABC$. If
\[
DE\parallel BC,
\]
then
\[
\frac{AD}{AB}
=
\frac{AE}{AC}
=
\frac{DE}{BC}.
\]
\end{theorembox}

\begin{center}
\begin{tikzpicture}[scale=0.9]
    \coordinate (A) at (0,3.6);
    \coordinate (B) at (-3.2,0);
    \coordinate (C) at (3.8,0);
    \coordinate (D) at (-1.6,1.8);
    \coordinate (E) at (1.9,1.8);
    \draw[subset] (A) -- (B) -- (C) -- cycle;
    \draw[very thick] (D) -- (E);
    \fill[geometricSubsetColor] (A) circle (2pt) node[above] {$A$};
    \fill[geometricSubsetColor] (B) circle (2pt) node[below left] {$B$};
    \fill[geometricSubsetColor] (C) circle (2pt) node[below right] {$C$};
    \fill[geometricSubsetColor] (D) circle (2pt) node[left] {$D$};
    \fill[geometricSubsetColor] (E) circle (2pt) node[right] {$E$};
    \node at (0.15,2.05) {$DE\parallel BC$};
\end{tikzpicture}
\end{center}

The smaller triangle $ADE$ and the larger triangle $ABC$ have the same three
angles. Their corresponding sides differ by one common scale factor. This is
the basic phenomenon of similarity.

\begin{definitionbox}
Two triangles are similar when their corresponding angles are equal. In that
case their corresponding side lengths are proportional.
\end{definitionbox}

Similarity lets us change the size of a figure while preserving its shape. It
is the reason ratios of corresponding lengths can depend only on angles.

\subsection{Sine and Cosine Before Coordinates}

Take a right triangle containing an acute angle $\alpha$. Let $h$ denote the
hypotenuse, let $a$ be the leg adjacent to $\alpha$, and let $b$ be the leg
opposite $\alpha$.

\begin{center}
\begin{tikzpicture}[scale=0.95]
    \coordinate (A) at (0,0);
    \coordinate (B) at (4.8,0);
    \coordinate (C) at (4.8,2.9);
    \coordinate (D) at (3.0,0);
    \coordinate (E) at (3.0,1.8125);

    % Large right triangle
    \draw[subset] (A) -- (B) -- (C) -- cycle;
    \draw (B) ++(-0.3,0) -- ++(0,0.3) -- ++(0.3,0);

    % Smaller similar triangle cut out by a parallel segment
    \draw[orange!85!black,line width=1pt] (D) -- (E);
    \draw (D) ++(-0.22,0) -- ++(0,0.22) -- ++(0.22,0);
    \draw[orange!85!black,->] (3.13,0.63) -- (3.13,1.03);
    \draw[orange!85!black,->] (4.93,1.05) -- (4.93,1.45);

    % Common angle
    \draw (0.85,0) arc[start angle=0,end angle=31.1,radius=0.85];
    \node at (1.08,0.28) {$\alpha$};

    % Side labels for the large triangle
    \node[below] at (3.9,0) {$a$};
    \node[right] at (4.8,2.05) {$b$};
    \node[above left] at (3.9,2.35) {$h$};

    % Side labels for the smaller triangle
    \node[below] at (1.5,0) {$a'$};
    \node[left,orange!85!black] at (3.0,0.92) {$b'$};
    \node[above left] at (1.55,0.98) {$h'$};

    \node[orange!85!black,align=left,anchor=west] at (5.35,1.45)
        {$DE\parallel BC$\\[-1pt]\small Thales' theorem};
\end{tikzpicture}

\[
\frac{a'}{a}=\frac{b'}{b}=\frac{h'}{h},
\qquad
\frac{a'}{h'}=\frac{a}{h},
\qquad
\frac{b'}{h'}=\frac{b}{h}.
\]
\end{center}

For an acute geometric angle $\alpha$, define
\[
\cos\alpha=\frac{a}{h},
\qquad
\sin\alpha=\frac{b}{h}.
\]
These ratios are independent of the chosen right triangle because all right
triangles containing the same acute angle are similar.

For angles occurring in an arbitrary Euclidean triangle, we also need the right
and obtuse cases. We set
\[
\cos(\text{right angle})=0.
\]
If $\alpha$ is obtuse and $\beta$ is its acute supplement, so that together they
form a straight angle, define
\[
\cos\alpha=-\cos\beta.
\]
Thus the geometric cosine is defined for every angle between $0$ and a straight
angle. For acute angles it is positive, for a right angle it is zero, and for
obtuse angles it is negative.

The essential role of Thales' theorem and similarity is that these values depend
only on the angle, not on the size of the triangle used to represent it.

For an acute angle, the Pythagorean theorem gives
\[
a^2+b^2=h^2.
\]
Dividing by $h^2$ yields
\[
\left(\frac{a}{h}\right)^2+
\left(\frac{b}{h}\right)^2=1,
\]
and hence
\[
\cos^2\alpha+\sin^2\alpha=1.
\]
Thus the fundamental trigonometric identity is a direct reformulation of the
Pythagorean theorem.
\subsection{Angles and Their Numerical Measure}

An angle exists geometrically before it is assigned a number. We can compare,
copy, add, and bisect angles without degrees or radians. A numerical angle
measure requires an additional choice of unit.

Degrees divide one full turn into $360$ equal parts. Radians use the geometry
of a circle. If an angle at the centre of a circle cuts off an arc of length
$s$ on a circle of radius $r$, its radian measure is
\[
\theta=\frac{s}{r}.
\]
The ratio does not depend on the size of the circle: scaling the circle by a
factor scales both $s$ and $r$ by that factor. Again, similarity makes the
quantity well defined.

\begin{remarkbox}
At this point no vector has been introduced. We have only Euclidean objects,
postulates, constructions, similarity, trigonometric ratios, and numerical
angle measure. Vectors will appear only in the next chapter, after the
Cartesian coordinate system has been constructed.
\end{remarkbox}

The next step is to choose two perpendicular lines, an origin, directions, and
a unit length. Those choices will convert the Euclidean plane into a Cartesian
coordinate plane.


\noindent\textbf{What the Cartesian construction uses.} The coordinate plane needs perpendicular axes, an origin, chosen directions, a common unit, and orthogonal projection. The detailed Euclidean reference explains the geometric background of those ingredients; the Cartesian section shows how they become a numerical language.
